\subsection{An independently adapted diffusion baseline}\label{sec:edm}
We adapt the official EGNN/EDM backbone \citep{edm} to coordinate generation with atom categories fixed as conditions. Positions follow its variance-preserving schedule. Starting from the GEOM-pretrained checkpoint, each run uses the same 3,000 OMol25 references as the source-model comparison, 30,000 updates, and the original $10^{-4}$ learning rate. We evaluate the 24 new compositions using 128-call and full 1,001-call samplers, including the final observation step.

\begin{table}[t]
\centering
\caption{Raw graph validity of independently adapted conditional EDM, with 768 attempts per run. The different-backbone harmonic and force-updated models reach 52.73\% and 53.45\% at 128 calls (Table~\ref{tab:fresh-graph}).}
\begin{tabular}{lrrr}
\toprule
Sampler & Run 1 & Run 2 & Pooled rate\\
\midrule
EDM, 128 calls &103&100&13.22\%\\
EDM, 1,001 calls &93&106&12.96\%\\
\bottomrule
\end{tabular}\label{tab:edm}
\end{table}

The task and adaptation data match, but architectures and pretraining differ: EDM has 2.38 million parameters and one pass per update, versus 5.90 million and two passes for LatentSpring. EDM fine-tuning takes 481.51 and 463.61 seconds, compared with approximately 570--590 seconds for the source models. Generating 768 outputs takes 48.54/47.94 seconds at 128 calls and 377.65/374.12 seconds at 1,001 calls. Thus call counts do not equate runtime. Composition-overlap checks cover the processed OMol25 corpora only. The matched Gaussian--harmonic comparisons provide the controlled source test.
